% Copyright (c) 2026 Christian Reitwiessner.
% Authors: Christian Reitwiessner
% Build: /tmp/andreev-pdf-tools/tectonic --keep-logs paper.tex
\documentclass[11pt]{article}
\usepackage[T1]{fontenc}
\usepackage{lmodern,amsmath,amssymb,amsthm,booktabs,microtype,enumitem}
\usepackage[margin=1in]{geometry}
\usepackage[colorlinks=true,allcolors=blue]{hyperref}
\newtheorem{theorem}{Theorem}[section]
\newtheorem{lemma}[theorem]{Lemma}
\newtheorem{proposition}[theorem]{Proposition}
\newtheorem{corollary}[theorem]{Corollary}
\newcommand{\bits}{\{0,1\}}
\newcommand{\ind}[1]{[\,#1\,]}
\newcommand{\THR}{\mathsf{THR}}
\newcommand{\smalltitle}[1]{\paragraph{#1}}
\setlength{\emergencystretch}{1em}
\setlist[itemize]{itemsep=1pt,topsep=3pt}
\setlist[enumerate]{itemsep=1pt,topsep=3pt}
\title{Exact and Approximate Threshold Circuits\\for Andreev's Function}
\author{Christian Reitwiessner}
\date{November 2026}
\begin{document}
\maketitle
% Copyright (c) 2026 Christian Reitwiessner.
% Authors: Christian Reitwiessner
\begin{abstract}
\marginpar{\raggedright\scriptsize\textbf{TODO:} The abstract still needs considerable work.}
Andreev's function reads an independent input table at an address formed
by parities of disjoint blocks. We give quadratic-size depth-two threshold
circuits for this function, improving the $O(n^3/\log n)$ upper bound
stated by Kane and Williams. A refinement uses $(5/16+o(1))n^2$ gates.
Under uniform inputs, explicit circuits achieve error $\delta$ with
$O(n^{3/2}\sqrt{\log(1/\delta)})$ gates, for $0<\delta\le1/2$.
For fixed $0<\delta<1/2$, the exponent matches their lower bound up to logarithms.
The approximation proof uses a symmetry of random Walsh coefficients:
odd errors in distinct coefficients are orthogonal even though those
coefficients are not independent. Both upper bounds use gates whose
address weights have one sign within each block. For this class we prove
an exact lower bound of $\Omega(n^2/\log^3 n)$ by comparing the information
preserved by a restriction with the number of restrictions one gate can
cross. We also give a size--advantage tradeoff and record a separate
exponential obstruction for full-support, equal-magnitude address weights.
\end{abstract}

\section{Introduction}
A depth-two threshold circuit has threshold gates on the raw input bits
and one threshold gate combining their outputs. We allow arbitrary real
weights and direct input wires to the output gate. The size is the number
of gates; intermediate counts below count bottom gates and add one output
gate. A wire is a nonzero edge after coefficients on the same edge are
aggregated.

Fix $k,m\ge1$ and $N=2^k$. The input consists of an independent $N$-bit
table $x=(x_c)_{c\in\bits^k}$ and $k$ address blocks $a_i\in\bits^m$.
The ordinary Andreev function is
$$
b_i=|a_i|\bmod2,\qquad A_{k,m}(x,a)=x_b,\qquad n=N+km.
$$
For asymptotic statements we take $m=\lfloor N/k\rfloor$, giving $n\sim2N$.
We also discuss exact balance $km=N$ when $k\mid2^k$. All finite-parameter
statements specify their own hypotheses.

\paragraph{The main comparison.}
Section~\ref{sec:identities} proves $O(n^2)$ exact size by a short
compilation argument. Rather than compute the address parities at an
extra layer, it writes lookup as a sum of functions of raw affine forms.
A threshold staircase compiles each such function directly into bottom
gates. Appendix~\ref{sec:exact} improves the leading constant to $5/16$;
none of the other proofs uses this refinement.

Section~\ref{sec:approx} gives explicit circuits of size
$O(n^{3/2}\sqrt{\log(1/\delta)})$ at uniform-input error $\delta$.
Symmetric truncation makes the error of each Fourier term odd in its
table coefficient. A table symmetry then cancels all cross-error moments,
so only individual concentration estimates are needed. A partial Fourier
sum gives agreement $1/2+\varepsilon$ using
$O(\varepsilon^2n^{3/2}\sqrt{\log(1/\varepsilon)})$ gates for
$N^{-1/2}\le\varepsilon\le1/8$.

Call a bottom gate \emph{sign-coherent} if its address weights in each
block are all nonnegative or all nonpositive; orientations may differ
between blocks and gates. Both constructions above are sign-coherent.
Section~\ref{sec:coherent} proves that exact computation in this class
requires $\Omega(n^2/\log^3 n)$ gates, with no restriction on table or
top weights. Thus exact and fixed-error complexity in this class have
exponents $2$ and $3/2$, respectively, up to logarithms.

\paragraph{Relation to earlier work.}
Kane and Williams~\cite[Theorems~1.1--1.2]{KW} prove lower bounds
$\Omega(\varepsilon^3n^{3/2}/\log^3 n)$ gates and
$\Omega(\varepsilon^3n^{5/2}/\log^{7/2}n)$ wires for agreement
$1/2+\varepsilon$, when $\varepsilon\gg\sqrt{\log n/n}$.
Their Section~5.2 states an $O(n^3/\log n)$ threshold upper bound and
a quadratic upper bound in the different $\mathsf{MOD}_3\circ\mathsf{MOD}_2$
basis. They expand lookup tests into parities and simulate those parities
separately. Our quadratic bound groups terms by address character first.
Their Section~4 already uses central-window staircases to approximate
parity, and their Section~6 notes the $\Omega(N^{-1/2})$ correlation with
majority of the table. Here the extra ingredients are odd-error
orthogonality and its use for partial lookup votes.

Threshold staircases are classical; see Hajnal et al.~\cite[Proposition~2.1]{HMPST}
for symmetric functions and Hansen--Podolskii~\cite[Theorem~5]{HP} for a
related representation using bounded integer forms. Our lower bound
combines a restriction information count with the single-crossing
principle used in threshold lower bounds~\cite{IPS}; it does not assume
bounded integer weights. Polynomial-overhead monotone-bottom simulations
\cite[Lemma~6]{HP} do not preserve the quantitative size bounds here.
Generalized Andreev functions, such as the encoded-table, extractor-address
function used by Chen, Santhanam and Srinivasan~\cite[Section~2.4]{CSS}, have
different inputs and are not the function considered here. More recent
$n^{5/2-\eta}$ gate lower bounds for functions in $E^{\mathsf{NP}}$
\cite{CTW} concern other functions; our remaining exponent gap is specific
to ordinary Andreev lookup.

For completeness, Appendix~\ref{sec:fullsign} records an elementary
full-support, equal-magnitude obstruction using the classical parity-order
argument of Minsky and Papert~\cite[Section~3.1]{MP}. It is a separate
restricted-class result, not an unrestricted lower-bound improvement.
We make no exhaustive priority claim for the elementary ingredients.

\section{Compiling functions of affine forms}\label{sec:identities}
The following elementary rule lets us reason about a numerical score
without introducing hidden circuit layers. Write $[P]$ for the Boolean
indicator of a statement $P$.

\begin{lemma}[Affine-form compilation]\label{lem:ladder}
Let $L$ be an integer-valued affine form in the raw inputs, with range
contained in $\{\ell,\ldots,h\}$. For any $g$ on this interval,
$$
g(L)=g(\ell)+\sum_{\theta=\ell+1}^{h}
\bigl(g(\theta)-g(\theta-1)\bigr)[L\ge\theta].
$$
Consequently a score consisting of an affine term plus a sum of such
$g(L)$ terms is realized by a weighted sum of bottom threshold gates.
Each nonzero finite difference costs at most one gate; repeated gates
can be shared.
\end{lemma}
\begin{proof}
The sum telescopes because the active comparisons are exactly those
with $\theta\le L$. All comparisons are on raw affine forms, and the
output threshold absorbs constants and affine terms.
\end{proof}

For the upper bounds, set
$$
X_S=\sum_{c\in\bits^k}(-1)^{S\cdot c+x_c},\qquad
t_S=\sum_{i\in S}|a_i|,\qquad S\subseteq[k].
$$
The coefficient $X_S$ is an even integer in $[-N,N]$.
Both $X_S$ and $t_S$ are affine in raw inputs, since
$(-1)^{x_c}=1-2x_c$. Character orthogonality gives the lookup identity
\begin{equation}\label{eq:lookup}
\sum_{S\subseteq[k]}X_S(-1)^{t_S}=N(-1)^{x_b}.
\end{equation}
Indeed, after exchanging sums, the coefficient of $(-1)^{x_c}$ is
$\prod_{i=1}^k(1+(-1)^{c_i+b_i})$, which is $N$ for $c=b$ and zero otherwise.
This is Fourier inversion for the table, with the address determining
which coordinate to recover. In linear-algebra notation, put
$v_c=(-1)^{x_c}$ and $H_{S,c}=(-1)^{S\cdot c}$. Then
$X=Hv$, $H^{\mathsf T}H=NI$, and $v=N^{-1}H^{\mathsf T}X$;
\eqref{eq:lookup} recovers the $b$th coordinate. Thus $X_S/N$ are the
coordinates of the truth table in the parity basis, rather than the point
basis; see also \cite[Chapter~1]{OD}.

\begin{theorem}[A quadratic exact circuit]\label{thm:simple}
For all $k,m\ge1$, $A_{k,m}$ has a sign-coherent circuit with at most
$2N(N-1)+kmN$ bottom gates and one output gate, using polynomially
bounded integer weights. In the balanced family this is
$(3/4+o(1))n^2$ gates.
\end{theorem}
\begin{proof}
For integers $u$, put $\varphi(u)=(-1)^u u$. If $X$ is even, then
$$X(-1)^t=\varphi(X+t)-\varphi(t).$$
Thus every nonempty summand of \eqref{eq:lookup} is a difference of two
functions of raw affine forms. Their range lengths are $2N+m|S|$ and
$m|S|$. The empty summand $X_\varnothing$ is already affine.
Lemma~\ref{lem:ladder} and $\sum_S|S|=kN/2$ give the count.
The score has values $\pm N$, so its negative is nonnegative exactly
when $x_b=1$. Address coefficients at every bottom gate are $0$ or $1$.
Bottom coefficients are bounded by $2$, thresholds by $O(n)$, and finite
differences of $\varphi$ by $O(n)$; aggregation preserves polynomial bounds.
\end{proof}

Appendix~\ref{sec:exact} reduces the constant by replacing equality of
scores with equality of signs. The main structural fact is already
visible here: parity times an even affine coefficient is itself a
difference of two functions of affine forms.

\section{Approximation from orthogonal odd errors}\label{sec:approx}
Write $\varphi(u)=(-1)^u u$ for integers $u$. All probabilities in this
section use independent uniform table and address bits. Balanced means $m=\lfloor N/k\rfloor$ or $km=N$.
Logarithms are natural unless a base is displayed.

\begin{theorem}\label{thm:avgupper}
In the balanced case there are explicit sign-coherent depth-two circuits
with the following guarantees:
\begin{enumerate}
\item Error at most $\delta$ with
 $O(n^{3/2}\sqrt{\log(1/\delta)})$ gates, for $0<\delta\le1/2$.
\item Agreement at least $1/2+\varepsilon$ with
 $O(\varepsilon^2n^{3/2}\sqrt{\log(1/\varepsilon)})$ gates and
 $O(\varepsilon^2n^{5/2}\sqrt{\log(1/\varepsilon)})$ wires, for
 $N^{-1/2}\le\varepsilon\le1/8$.
\end{enumerate}
For $1/8<\varepsilon<1/2$, Part~1 applies with $\delta=1/2-\varepsilon$.
\end{theorem}

The key to summing truncation errors is a symmetry stronger than
uncorrelatedness of the coefficients themselves.
\begin{lemma}[Odd-error orthogonality]\label{lem:odd}
For a uniform independent table $x$, distinct $S,T\subseteq[k]$,
an odd function $g:\mathbb Z\to\mathbb R$, and any $h:\mathbb Z\to\mathbb R$,
$$\mathbb E_x[g(X_S)h(X_T)]=0.$$
The functions may depend on any fixed address input.
\end{lemma}
\begin{proof}
Choose $u\in\bits^k$ with $(S\mathbin\triangle T)\cdot u$ odd.
For $v_c=(-1)^{x_c}$, the signed permutation
$v'_c=(-1)^{T\cdot u}v_{c\oplus u}$ preserves the distribution of $v$.
It changes $X_S$ to $-X_S$ and leaves $X_T$ fixed. The expectation is
therefore its own negative.
\end{proof}
Walsh coefficients are not independent. The lemma uses their pairwise
sign symmetry and applies after conditioning on the address input.

\begin{proof}[Proof of Theorem~\ref{thm:avgupper}]
\emph{Partial lookup votes.}
Choose $0\le r\le k$ and put $q=2^r$. Orthogonality gives
$$
P_*=2\sum_{S\subseteq[r]}X_S(-1)^{t_S}
 =2q\sum_{c:c_{[r]}=b_{[r]}}(-1)^{x_c}.
$$
This is a majority vote of $L=N/q$ table bits including $x_b$.
For $q=N$ it is exact. Otherwise $L$ is even, and with ties outputting
zero its advantage is $p(L)/2$, where
$$
p(L)=2^{-L}\binom L{L/2}\ge(2L)^{-1/2}.
$$
To see the agreement formula, condition on the chosen bit; the other
$L-1$ signs have sum $-1$ with probability $p(L)$, which gives excess
agreement $p(L)/2$. The inequality follows from $p(2)=1/2$ and
$p(L+2)/p(L)=(L+1)/(L+2)\ge\sqrt{L/(L+2)}$.
The vote is $[-P_*\ge2q]$ for $q<N$ and $[-P_*\ge0]$ for $q=N$.

\emph{Symmetric truncation.}
For nonempty $S\subseteq[r]$, center the count by setting
$$D_S=t_S-\lfloor m|S|/2\rfloor,\qquad
\eta_S=(-1)^{\lfloor m|S|/2\rfloor}.$$
Since $X_S$ is even,
$$
\eta_S\bigl(\varphi(X_S+D_S)+\varphi(X_S-D_S)\bigr)
 =2X_S(-1)^{t_S}.
$$
For an integer $R\ge1$ let $f_R(u)=\varphi(u)[|u|\le R]$ and define
$$
P_R=2X_\varnothing+
\sum_{\varnothing\ne S\subseteq[r]}
\eta_S\bigl(f_R(X_S+D_S)+f_R(X_S-D_S)\bigr).
$$
The identity
$$
f_R(U)=\sum_{\ell=-R}^{R+1}
\bigl(f_R(\ell)-f_R(\ell-1)\bigr)[U\ge\ell]
$$
uses at most $2R+2$ thresholds for every integer $U$.
Thus $P_R$ uses at most $4(q-1)(R+1)$ bottom gates. The address weights
in $X_S+D_S$ are $+1$ on selected blocks, and those in $X_S-D_S$ are
$-1$ there; other blocks have weight zero. These gates are sign-coherent.

\emph{The mean-square estimate.}
Let $e_S$ be a truncated summand minus its exact value.
Both $\varphi$ and $f_R$ are odd. Hence, for every fixed $a$, $e_S$ is
an odd function of $X_S$: replacing $X_S$ by $-X_S$ interchanges and
negates the two terms. Lemma~\ref{lem:odd} gives
\begin{equation}\label{eq:odd-l2}
\mathbb E_{x,a}(P_R-P_*)^2=\sum_{S\ne\varnothing}\mathbb E_{x,a}e_S^2.
\end{equation}
This is why the two reflected forms are used; controlling each term
separately and taking a union bound would lose a logarithmic factor.

For $U=X_S\pm D_S$, the mean has magnitude at most $1/2$ and its centered
version is subgaussian with variance proxy $v_S=N+m|S|/4$.
Indeed, the table contributes $N$ independent signs and the centered
address count contributes $m|S|$ independent signs of magnitude $1/2$.
Writing $V=U-\mathbb EU$, Hoeffding's bound is
$\Pr(|V|>t)\le2e^{-t^2/(2v_S)}$.
Since $|U|>R$ implies $|V|>R/2$ and $U^2\le2V^2+1/2$, tail integration gives
$$
\mathbb E[U^2[|U|>R]]
 \le(R^2+8v_S+1)e^{-R^2/(8v_S)}
 \le16v_Se^{-R^2/(16v_S)}.
$$
For the last inequality use $v_S\ge2$ and $1+t\le e^t$.
As $|\varphi(U)-f_R(U)|=|U|[|U|>R]$, the inequality
$(u+v)^2\le2u^2+2v^2$ and \eqref{eq:odd-l2}, with $v=N+rm/4$, give
$$
\mathbb E(P_R-P_*)^2\le64(q-1)v e^{-R^2/(16v)}.
$$

\emph{From scores to agreement.}
When $q<N$, $P_*$ takes values in $4q\mathbb Z$, so the vote threshold
$2q$ has distance at least $2q$ from every attainable score.
For $q=N$, the scores are $\pm2N$ and the threshold is zero.
In either case the circuit using $P_R$ can disagree with the vote only
if $|P_R-P_*|\ge2q$. Its disagreement probability is therefore at most
$$
\frac{16v}{q}e^{-R^2/(16v)}.
$$
Choosing
$$R=\left\lceil4\sqrt{v\log\frac{16v}{q\delta}}\right\rceil$$
makes this at most $\delta$. This specifies one deterministic circuit
directly, without a choice of random masks. One may cap $R$ at $N+km$,
since truncation is then exact on all attainable forms.

For Part~1 set $q=N$. Balance gives $v=O(N)$ and
$R=O(\sqrt{N\log(1/\delta)})$. For Part~2 let $q$ be the least power of
two at least $32\varepsilon^2N$ and set $\delta=\varepsilon$.
Then $q<64\varepsilon^2N$, $q\le N$, and the vote's advantage is at least
$\sqrt{q/(8N)}\ge2\varepsilon$; if $q=N$ its advantage is $1/2$.
Also $\log(16v/(q\delta))=O(\log(1/\varepsilon))$, proving the gate bound.
A circuit with $G$ bottoms has at most $nG+G+n$ support wires, including
direct top wires; the stated parameter range absorbs the additive terms.
\end{proof}

\section{Restriction demand and crossing capacity}\label{sec:coherent}
The lower bound compares two quantities. Every restriction in a family
must preserve enough gates to describe the table; each coherent gate
can remain nonconstant on only a small fraction of that family.
If the average demand is $D$ and the per-gate fraction is at most $p$,
double counting gives $Gp\ge D$. We make both estimates explicit.

\begin{lemma}[Information demand under a restriction]\label{lem:entropy}
Consider a fixed depth-two threshold circuit with $G$ bottoms, input table
$x\in\bits^N$, and $h\ge1$ additional input bits $z$. Suppose that the function
of $z$ computed for each $x$ determines $x$ uniquely. If $D(x)$ bottoms
are nonconstant as functions of $z$, then for uniform $x$,
$$
\mathbb E D(x)\ge
\frac{\max\{0,N-h-1-\log_2(G+1)\}}{h+\log_2(G+1)}.
$$
All weights may be arbitrary real numbers, and the top gate may read
both $x$ and $z$ directly.
\end{lemma}
\begin{proof}
Only the threshold of each bottom gate changes with $x$; its coefficients
on $z$ are fixed. A fixed normal has at most $2^h-1$ nonconstant threshold
functions on $\bits^h$. Specify the set of $D$ nonconstant gates and one
such function for each. The remaining constant outputs and direct table
inputs are absorbed into the top threshold. Its other coefficients are
fixed, so it has at most $2^h+1$ possible functions after the specification.
This description determines the output function and hence $x$.
Using base-two entropy,
\begin{align*}
N=H(x)&\le H(D)+\mathbb E\log_2\left[
 \binom GD(2^h-1)^D(2^h+1)\right]\\
&\le\log_2(G+1)+\mathbb E D\bigl(\log_2(G+1)+h\bigr)+h+1.
\end{align*}
Rearrange and use $D\ge0$.
\end{proof}

\begin{theorem}\label{thm:coherent}
For $N=2^k$, every exact sign-coherent circuit for $A_{k,m}$ with $G$
bottom gates satisfies
$$
G\ge\frac{D_0}{1-(1-1/m)^k}\ge\frac m kD_0,
\qquad
D_0=\frac{\max\{0,N-k-1-\log_2(G+1)\}}{k+\log_2(G+1)}.
$$
At floor or exact balance this is $\Omega(n^2/\log^3 n)$.
\end{theorem}
\begin{proof}
\emph{Each cube reveals the table.}
Fix an order within each block. Let $a(s)$ have ones in exactly the first
$s_i$ positions of block $i$, for $s\in\{0,\ldots,m\}^k$.
For each cube $q\in\{0,\ldots,m-1\}^k$, the inputs $a(q+z)$,
$z\in\bits^k$, vary one raw bit per block and satisfy
$$A_{k,m}(x,a(q+z))=x_{(q\bmod2)\oplus z}.$$
Thus every cube reveals all $N$ table bits. Applying
Lemma~\ref{lem:entropy} with $h=k$ shows that the expected number of
bottom gates splitting any fixed cube is at least $D_0$.

\emph{One gate crosses few cubes.}
Fix a table and a gate. Choose a direction $\epsilon_i\in\{-1,1\}$ in
which its block score is nondecreasing. The gate's score on
$a(s)$ is monotone when each $s_i$ moves in direction $\epsilon_i$;
unequal magnitudes and zero weights cause no difficulty.
On cube $q$ its minimum and maximum occur at
$$q^-=q+(\mathbf1-\epsilon)/2,\qquad
q^+=q+(\mathbf1+\epsilon)/2.$$
Writing $F$ for the raw affine score and $\theta$ for the threshold, the
gate splits the cube exactly when $F(q^-)<\theta\le F(q^+)$.

Partition the cubes into lines in direction $\epsilon$.
The maximum vertex of one cube is the minimum vertex of the next, so
the half-open intervals $(F(q^-),F(q^+)]$ on a line are disjoint.
A gate splits at most one cube per line. Each line starts at a cube
$q$ for which $q-\epsilon$ lies outside the grid. Exactly $(m-1)^k$ cubes
have a predecessor, giving $m^k-(m-1)^k$ lines.

\emph{Counting incidences.}
Sum the demand over all $m^k$ cubes and the capacity over all gates:
$$m^kD_0\le G\bigl(m^k-(m-1)^k\bigr).$$
Use $1-(1-1/m)^k\le k/m$ for the second inequality.
If $G\ge N^2$ the asymptotic conclusion is immediate. Otherwise integrality
gives $G+1\le N^2$, and $D_0\ge(N-3k-1)/(3k)$ for large $k$.
Finally $m=\Theta(N/k)$ and $n=\Theta(N)$.
\end{proof}

\paragraph{Why this is an exact lower bound.}
The cube restrictions sample all block counts, not uniform raw addresses.
An error guarantee under uniform inputs need not survive this distribution.
The theorem therefore coexists with the sign-coherent approximate circuits
of Section~\ref{sec:approx}. To beat the exact exponent $2$ by a fixed
amount, a construction must use mixed signs within at least one block.

\section{Further questions}
The unrestricted exact gap remains between exponents $3/2$ and $2$.
Can mixed signs inside address blocks yield $n^{2-\eta}$ gates for some
fixed $\eta>0$? Can exact wire size improve on the $O(n^3)$ consequence
of our gate bound, towards the known lower exponent $5/2$? A gate upper
bound transfers to wires by multiplying by $n$; a stronger wire lower
bound needs separate support information.

For approximation at fixed error below $1/2$, the exponent is settled up to logarithms,
but the dependence on error is not. Our bound reaches quadratic size
when $\log(1/\delta)=\Theta(n)$, and the small-advantage upper bound has
an $\varepsilon^2$ dependence versus $\varepsilon^3$ in the known lower
bound. Closing either gap would clarify the transition from weak
prediction to exact computation. A further question is which structured
or biased table distributions retain enough of the sign symmetry in
Lemma~\ref{lem:odd} for the same approximation argument to work.

The research and exposition were developed with AI assistance. Exact checking scripts accompany the source: \texttt{verify\_results.py}
checks the refined construction, \texttt{approximation\_checks.py} checks
the explicit truncation and its orthogonality, and
\texttt{restriction\_checks.py} checks the lower-bound restrictions.
They supplement the proofs. An earlier quadratic construction is
formalized in Lean in \href{https://github.com/cslib-community/cslib-frontier/pull/2813}{cslib-frontier};
the other results here are not claimed to be formalized.

\appendix
\section{The sharper exact construction}\label{sec:exact}
This appendix improves only the constant in Theorem~\ref{thm:simple}.
One staircase per character replaces the two in the basic construction.
The price is a quadratic correction, whose sum collapses to one global
square and $k$ block squares. Corrections at the zero and singleton
addresses then enforce the desired sign. All auxiliary notation below
is local to this appendix.
\begin{theorem}\label{thm:exact}
Put $\alpha=\lfloor(m-1)/2\rfloor$. There is an exact depth-two threshold
circuit for $A$ with at most
\begin{equation}\label{eq:bestcount}
\frac{N(N-1)}2+\frac{kN(m+\alpha)}2
 +2(N-1)+2\alpha k+3km+\left\lceil\frac{km}{2}\right\rceil
\end{equation}
bottom gates and one top gate. Its weights are integers of magnitude
polynomial in $n$. At exact balance its total size is
$$
\frac5{16}n^2-\frac n4\log_2(n/2)+O(n).
$$
For $m=\lfloor N/k\rfloor$, its size is $(5/16+o(1))n^2$.
\end{theorem}

\begin{proof}
\emph{Forms used only in this proof.}
Write $s_i=|a_i|$, $\sigma=\sum_i s_i$, and let $e_i$ be the address with
only coordinate $i$ equal to one; $x_0$ means the all-zero address entry. Put
$$
z=\sum_{c\ne0}x_c,\quad E=\sum_i x_{e_i},\quad
w=2z+2\alpha E+\sigma,\quad R=2(N-1)+2\alpha k+km,\quad C=N(R+4).
$$
Thus $0\le w\le R$. For each nonempty $S\subseteq[k]$, define
$$
W_S=\sum_{c:S\cdot c\text{ odd}}x_c+\alpha\sum_{i\in S}x_{e_i},\quad
\kappa_S=-4(2\alpha+1)(k-2|S|),
$$
and let $r_S(u)=-8\lceil u/2\rceil^2+\kappa_S\lceil u/2\rceil$.
All of $W_S,t_S,w,s_i$ are affine forms in the original inputs.

\emph{The score.} Consider
\begin{align}\label{eq:bestscore}
P={}&2N+\sum_{S\ne\varnothing}r_S(2W_S+t_S)
 +\frac N2\left(w^2+\sum_i s_i^2\right)\notag\\
 &+4Nz-2N(\alpha^2+1)E-2\alpha N\sigma
 +2N(1+\alpha)\sum_i x_{e_i}s_i-Cx_0\ind{b=0}.
\end{align}
We first verify its sign, then realize its non-affine terms with bottom gates.
The following identities follow by expanding the squares and summing
characters; all sums over $S$ are over nonempty sets:
\begin{align*}
\sum_S4W_S(-1)^{t_S}
 &=2N\ind{b=0}z-2N\ind{b\ne0}x_b
   +2\alpha N\sum_i x_{e_i}\bigl(\ind{b=0}-\ind{b=e_i}\bigr),\\
8\sum_S(W_S^2+W_St_S)
 &=\frac N2(w^2-\sigma^2)+2Nz+(4\alpha+2\alpha^2)NE
   +2N(1+\alpha)\sum_i x_{e_i}s_i,\\
8\sum_S\lceil t_S/2\rceil^2
 &=\frac N2\left(\sigma^2+\sum_i s_i^2\right)
   +N(\sigma+1)\ind{b\ne0}+N\sum_i s_i\ind{b=e_i}.
\end{align*}
To check the quadratic expansion, for distinct nonzero addresses $c,d$ use
$$\sum_S[S\cdot c\text{ odd}]=N/2,\qquad
\sum_S[S\cdot c\text{ odd}][S\cdot d\text{ odd}]=N/4.$$
For $c=d$ the second sum is $N/2$. These counts expand $W_S^2$ directly;
the other terms use the same orthogonality and
$\lceil t/2\rceil=(t+[t\text{ odd}])/2$.
Also
$$
4W(-1)^t=4W-8\lceil(2W+t)/2\rceil^2+8W^2+8Wt+8\lceil t/2\rceil^2.
$$
The shifts $\kappa_S$ change only coefficients of existing odd thresholds,
so the singleton-address correction adds no staircase gates. They supply
the address correction because
$$
\sum_S\kappa_S\bigl(\lceil t_S/2\rceil-t_S/2\bigr)
 =N(2\alpha+1)\sum_i\ind{b=e_i}.
$$
The remaining linear shifts satisfy
$$\sum_S\kappa_SW_S=2N(2\alpha+1)(1+\alpha)E,\qquad
\frac12\sum_S\kappa_St_S=N(2\alpha+1)\sigma.$$
These follow by writing $k-2|S|=\sum_i(-1)^{[i\in S]}$ and applying the
same orthogonality. Substitution in \eqref{eq:bestscore} gives
\begin{equation}\label{eq:bestclosed}
P=\begin{cases}
N(w+2)-Cx_0,&b=0,\\
2N(1+\alpha)(1-x_{e_i})-Ns_i,&b=e_i,\\
N(1-2x_b),&\text{otherwise}.
\end{cases}
\end{equation}
At $b=0$, use $0\le w\le R$ and $C=N(R+4)$. At $b=e_i$, the odd
integer $s_i$ lies between $1$ and $2\alpha+1$. In every case $P\ge N$
when $x_b=0$, and $P\le-N$ when $x_b=1$. The top gate outputs $\ind{-P\ge0}$.

\emph{Raw gates and their count.}
The staircase for $r_S(2W_S+t_S)$ uses only odd thresholds
$$
\ind{2W_S+t_S\ge\theta},\qquad
1\le\theta\le N+2\alpha|S|+|S|m,
$$
with coefficient $-8\theta+\kappa_S$. The two square terms use $R+km$
gates by Lemma~\ref{lem:ladder}. For $B=n+2$, each product $x_{e_i}s_i$
uses $m$ gates $\ind{s_i+Bx_{e_i}\ge B+\ell}$, $1\le\ell\le m$.
Finally, the identity
$$
x_0\ind{b=0}=x_0-x_0\sigma+
 \frac4N\sum_{S\ne\varnothing}x_0\lfloor t_S/2\rfloor
$$
implements the last correction with $km$ gates
$\ind{\sigma+Bx_0\ge B+\ell}$ and
$\sum_S\lfloor |S|m/2\rfloor$ gates
$\ind{t_S+Bx_0\ge B+2j}$. Its identity follows from
$\sum_S\ind{t_S\text{ odd}}=(N/2)\ind{b\ne0}$.
The even-level $\sigma$ gates coincide with the $S=[k]$ gates, saving
$\lfloor km/2\rfloor$ gates. Everything else in $P$ is an affine top term.

Since $\sum_{S\ne\varnothing}|S|=kN/2$, the main-ridge and last-correction
counts together are
$$
\sum_{S\ne\varnothing}\left(
 \left\lceil\frac{N+2\alpha|S|+|S|m}{2}\right\rceil
 +\left\lfloor\frac{|S|m}{2}\right\rfloor\right)
 =\frac{N(N-1)}2+\frac{\alpha kN}2+\frac{kmN}2.
$$
Adding the remaining gates gives \eqref{eq:bestcount}. Vanishing coefficients
or further coincidences can only improve it. At exact balance $m$ is even,
so $\alpha=m/2-1$; substitution gives the stated asymptotic bound.
For floor balance, the rounding changes the leading expression by
$O(Nk)=o(n^2)$. Bottom coefficients and thresholds are $O(n)$; the displayed
expansion has top coefficients bounded by $O(n^3)$.
\end{proof}



% Copyright (c) 2026 Christian Reitwiessner.
% Authors: Christian Reitwiessner
\section{A full-support equal-magnitude obstruction}\label{sec:fullsign}
Here the restriction demand is only one gate, but a gate can serve an
exponentially small fraction of the restrictions. The demand comes from
the classical parity-order obstruction~\cite[Section~3.1]{MP}.

\begin{lemma}[The missing-variable parity obstruction]\label{lem:paritydual}
Let $d\ge2$ and $f(z)=\prod_{i=1}^d z_i$ on $\{-1,1\}^d$.
If a score consisting of an affine term and a linear combination of
predicates strictly has the sign of $f$, then some predicate depends
on every variable.
\end{lemma}
\begin{proof}
A predicate missing coordinate $i$ has zero correlation with $f$:
summing over $z_i$ cancels. Affine terms also have zero correlation for
$d\ge2$. If every predicate missed a variable, the score $F$ would have
$\mathbb E[fF]=0$, whereas strict sign agreement gives
$\mathbb E[fF]=\mathbb E|F|>0$.
\end{proof}

\begin{proposition}\label{thm:fullsign}
Let $k\ge2$, $m\ge1$, and $d=m\lfloor k/2\rfloor\ge2$. Suppose every
bottom gate's address normal is either zero or $\rho\sigma$ with
$\rho>0$ and $\sigma\in\{-1,1\}^{km}$. An exact circuit for $A_{k,m}$
then has at least $2^d$ bottom gates. Table weights, thresholds, top
weights, and raw top wires are unrestricted real numbers.
At balance the bound is $2^{\Omega(n)}$.
\end{proposition}
\begin{proof}
Fix the table to parity of its first $r=\lfloor k/2\rfloor$ address
coordinates. If $k$ is odd, fix one of the other blocks to zero.
There remain $d=rm$ selected bits and $d$ ignored bits. In sign coordinates
$2a-1$, call these $z_i,w_i$. Up to a constant sign the target is
$f(z)=\prod_i z_i$.
We can make the top score strictly correct: for output $[F\ge0]$, add
a small positive constant before adjusting the orientation. This works
on the finite input domain because all negative-class scores were already
strictly negative.

For each $t\in\{-1,1\}^d$, identify $w_i=t_i z_i$.
These are literal or complemented-literal identifications, so raw top
terms remain affine. A nonzero-normal gate with paired signs
$\sigma_i,\tau_i$ acquires coefficients proportional to
$\sigma_i+t_i\tau_i$. It can depend on every $z_i$ only for the unique
choice $t_i=\sigma_i\tau_i$ for all $i$. A zero-normal gate is constant.
By Lemma~\ref{lem:paritydual}, every one of the $2^d$ restrictions needs
at least one gate depending on all live variables. Each original gate
can serve at most one restriction, proving the count.
\end{proof}

Use $X_S,t_S,\varphi$ from Section~\ref{sec:identities}.
Both qualifications are essential. Without full support,
Theorem~\ref{thm:simple} already gives quadratic size with address weights
in $\{0,1\}$. Without equal magnitudes, put
$$L_S=\sum_{i\in S,j}a_{i,j}+2\sum_{i\notin S,j}a_{i,j}.$$
This has full address support and the parity of $t_S$.
The identity $X_S(-1)^{t_S}=\varphi(X_S+L_S)-\varphi(L_S)$ and
Lemma~\ref{lem:ladder} give $O(n^2)$ gates, since each range has length
$O(n)$. Thus the proposition does not extend to general bounded integer
address weights, even when they have full support.

\begin{thebibliography}{9}\small
\bibitem{OD} R.~O'Donnell. \emph{Analysis of Boolean Functions}.
Cambridge University Press, 2014.
\href{https://www.cs.cmu.edu/~odonnell/papers/Analysis-of-Boolean-Functions-by-Ryan-ODonnell.pdf}{Author-hosted edition}.
\bibitem{CTW} L.~Chen, A.~Tal, and Y.~Wang.
Super-quadratic Lower Bounds for Depth-2 Linear Threshold Circuits.
\href{https://eccc.weizmann.ac.il/report/2026/039/}{ECCC TR26-039}, 2026.
\bibitem{KW} D.~M. Kane and R.~Williams.
Super-Linear Gate and Super-Quadratic Wire Lower Bounds for Depth-Two and
Depth-Three Threshold Circuits. STOC, 2016.
\href{https://arxiv.org/abs/1511.07860}{arXiv:1511.07860}.
\bibitem{HMPST} A.~Hajnal, W.~Maass, P.~Pudl\'ak, M.~Szegedy, and G.~Tur\'an.
Threshold circuits of bounded depth. J. Comput. Syst. Sci. 46(2):129--154, 1993.
\bibitem{HP} K.~A. Hansen and V.~V. Podolskii.
Polynomial threshold functions and Boolean threshold circuits.
Information and Computation 240:56--73, 2015.
Theorem and lemma numbering here follows
\href{https://eccc.weizmann.ac.il/report/2013/021/download/}{ECCC TR13-021}.
\bibitem{MP} M.~L. Minsky and S.~A. Papert.
\emph{Perceptrons: An Introduction to Computational Geometry}.
MIT Press, 1969; expanded edition, 1988. Section~3.1.
\bibitem{IPS} R.~Impagliazzo, R.~Paturi, and M.~Saks.
Size--depth tradeoffs for threshold circuits. SIAM J. Comput. 26(3):693--707, 1997.
\bibitem{CSS} R.~Chen, R.~Santhanam, and S.~Srinivasan.
Average-Case Lower Bounds and Satisfiability Algorithms for Small Threshold
Circuits. Theory of Computing 14(9):1--55, 2018.
\end{thebibliography}
\end{document}
